\documentclass[../../main.tex]{subfiles}
\begin{document}
\section{习题}\label{section:ch15-problems}

\begin{enumerate}
\item\label{exercise:ch15-01b} 设 $M$ 是一个光滑流形，它是两个可定向开子流形的并，且它们的交连通。证明 $M$ 可定向。利用这一点给出 $S^n$ 可定向的另一个证明。

\item\label{exercise:ch15-02} 设 $M,N$ 是带边或不带边的有向光滑流形，$F:M\to N$ 是一个局部微分同胚。证明：若 $M$ 连通，则 $F$ 要么保定向，要么反定向。

\item\label{exercise:ch15-03} 设 $n\geqslant1$，令 $\alpha:S^n\to S^n$ 为对径映射：$\alpha(x)=-x$。证明 $\alpha$ 保定向当且仅当 $n$ 为奇数。[提示：考虑由 $F(x)=-x$ 给出的映射 $F:\overline{\mathbb B}^n\to\overline{\mathbb B}^n$，并利用推论\ref{corollary:ch15-boundary-volume}。]（用于第~393、435 页。）

\item\label{exercise:ch15-04b} 令 $\Theta$ 是一个带边或不带边的有向光滑流形上的光滑流。证明：对每个 $t\in\mathbb{R}$，$\Theta_t$ 在其有定义处都保定向。（用于第~425 页。）

\item\label{exercise:ch15-05} 设 $M$ 是一个带边或不带边的光滑流形。证明 $TM$ 与 $T^*M$ 的全空间都可定向。

\item\label{exercise:ch15-06} 令 $U\subseteq\mathbb{R}^3$ 为开子集
\begin{align*}
\left\{(x,y,z):\left(\sqrt{x^2+y^2}-2\right)^2+z^2<1\right\}
\end{align*}
（例\ref{example:ch05-torus} 中旋转环面所围成的实心环）。由
\begin{align*}
F(u,v)=\bigl(\cos 2\pi u\,(2+\tanh v\cos\pi u),\sin 2\pi u\,(2+\tanh v\cos\pi u),\tanh v\sin\pi u\bigr)
\end{align*}
定义一个映射 $F:\mathbb{R}^2\to U$。
\begin{enumerate}[(a)]
\item 证明 $F$ 下降为从 $E$ 到 $U$ 的一个光滑嵌入，其中 $E$ 是例\ref{example:ch10-03} 中 Möbius 丛的全空间。
\item 令 $S=F(E)$。证明 $S$ 是 $U$ 的一个固有嵌入光滑子流形。
\item 证明沿 $S$ 不存在单位法向量场。
\item 证明 $S$ 在 $U$ 中没有整体定义函数。
\end{enumerate}

\item\label{exercise:ch15-07} 设 $M$ 是一个带边或不带边的有向 Riemann 流形，$S\subseteq M$ 是一个带边或不带边的有向光滑超曲面。证明：沿 $S$ 存在唯一一个光滑单位法向量场，它确定 $S$ 给定的定向。

\item\label{exercise:ch15-08b} 设 $M$ 是一个可定向 Riemann 流形，$S\subseteq M$ 是一个带边或不带边的浸入或嵌入子流形。证明下列命题。
\begin{enumerate}[(a)]
\item 若 $S$ 的法丛平凡，则 $S$ 可定向。
\item 若 $S$ 是一个可定向超曲面，则 $S$ 的法丛平凡。
\end{enumerate}

\item\label{exercise:ch15-09} 令 $S$ 是 $\mathbb{R}^3$ 中一个有向的、嵌入的、二维带边子流形，令 $C=\partial S$ 带诱导定向。由习题\ref{exercise:ch15-07}，$S$ 上存在唯一的光滑单位法向量场 $N$ 确定该定向。令 $T$ 为 $C$ 上定向的单位切向量场，$V$ 为 $C$ 沿线与 $T$ 正交、向内指向、且与 $S$ 相切的唯一单位向量场。证明：在每个 $p\in C$ 处，$(T_p,V_p,N_p)$ 是 $\mathbb{R}^3$ 的一组定向标准正交基。

\item\label{exercise:ch15-10b} \textbf{定向覆盖的特征性质：} 设 $M$ 是一个连通的、不可定向的、带边或不带边的光滑流形，$\widehat\pi:\widehat M\to M$ 是它的定向覆盖。证明：若 $X$ 是任意一个带边或不带边的、有向的光滑流形，$F:X\to M$ 是任意一个局部微分同胚，则存在唯一的保定向局部微分同胚 $\widehat F:X\to\widehat M$，使 $\widehat\pi\circ\widehat F=F$：
\begin{align*}
\begin{CD}
\widehat M\\
@V\widehat\pi VV\\
X @>F>> M.
\end{CD}
\end{align*}

\item\label{exercise:ch15-11} 证明定理\ref{theorem:ch15-covering-uniqueness}（定向覆盖的唯一性）。[提示：使用习题\ref{exercise:ch15-10b}。]

\item\label{exercise:ch15-12b} 证明 $\mathbb{R}$ 的每个反定向微分同胚都有不动点。

\item\label{exercise:ch15-13b} \textbf{光滑 $1$-流形的分类：} 设 $M$ 是一个连通的光滑 $1$-流形。按如下步骤证明 $M$ 微分同胚于 $\mathbb{R}$ 或 $S^1$ 之一：
\begin{enumerate}[(a)]
\item 首先处理 $M$ 可定向的情形：证明 $M$ 容许一个非零光滑向量场，并使用习题\ref{exercise:ch09-01}。
\item 现在设 $M$ 任意，证明 $M$ 可定向：证明它的万有覆盖流形微分同胚于 $\mathbb{R}$，并使用习题\ref{exercise:ch15-12b} 的结果。
\end{enumerate}
由此推出 $\mathbb{R}$ 与 $S^1$ 上的光滑结构在微分同胚意义下唯一。

\item\label{exercise:ch15-14b} \textbf{带边光滑 $1$-流形的分类：} 证明：每个边界非空的连通光滑 $1$-流形都微分同胚于 $[0,1]$ 或 $[0,\infty)$ 之一。[提示：使用倍流形。]

\item\label{exercise:ch15-15} 设 $M$ 是 $\mathbb{R}^n$ 中一个不可定向的嵌入超曲面，令 $NM$ 为它的法丛，投影为 $\pi_{NM}:NM\to M$。证明集合
\begin{align*}
W=\{(x,v)\in NM:|v|=1\}
\end{align*}
是 $NM$ 的一个嵌入子流形，且 $\pi_{NM}$ 在 $W$ 上的限制是一个光滑覆盖映射，同构于 $M$ 的定向覆盖。[提示：考虑由 $v\lrcorner(\mathrm{d}x^1\wedge\cdots\wedge\mathrm{d}x^n)$ 所确定的定向。]

\item\label{exercise:ch15-16b} 设 $E$ 是例\ref{example:ch15-mobius-bundle} 中 Möbius 丛的全空间。证明 $E$ 的定向覆盖微分同胚于柱面 $S^1\times\mathbb{R}$。
\end{enumerate}
\end{document}
